\section{Aguas profundas: intuición matemática, el ecosistema de Lean y la colaboración en la investigación}

Los capítulos anteriores desglosaron AI for Math en una pila tecnológica, pero lo verdaderamente interesante de la entrevista con Hong Letong (洪乐潼) es que ella no concibe las matemáticas como un simple «conjunto de problemas». En la investigación matemática hay estética, hay fracasos, hay dedicación a largo plazo y también hay reconocimiento de la comunidad. Si la IA solo sabe generar respuestas dentro de bancos de problemas ya conocidos, sigue muy lejos de la investigación matemática; solo cuando puede ayudar a los matemáticos a descubrir estructuras, formalizar demostraciones, completar los lemma faltantes y verificar contraejemplos entra de verdad en el proceso de investigación.

\subsection{Por qué las matemáticas no son una tarea común de NLP}

Las tareas comunes de NLP suelen admitir aproximaciones semánticas: un resumen puede redactarse de distintas maneras, una traducción puede tener varias expresiones equivalentes y las respuestas de un chat también admiten diferencias de estilo. Las demostraciones matemáticas son distintas: basta con omitir una condición clave para que toda la demostración pueda fallar. Los sistemas formales llevan este rigor al extremo y exigen que cada objeto, tipo, dependencia y paso de razonamiento sea válido.

\begin{importantbox}{La particularidad de las tareas matemáticas}
Las matemáticas tienen a la vez fuertes restricciones formales y una fuerte creatividad. Las restricciones formales provienen del proof assistant y de la verificación de teoremas; la creatividad proviene de las conjeturas, la estética, la elección de definiciones y el diseño de la ruta de demostración. AI for Math debe atender ambos extremos al mismo tiempo.
\end{importantbox}

\subsection{Por qué es importante el ecosistema de Lean}

El valor de Lean no reside solo en ser un verificador, sino en su ecosistema: la biblioteca de teoremas mathlib, las contribuciones de la comunidad, el estilo de formalización, la tradición de tactic, los materiales didácticos y la participación de los matemáticos. Para escribir demostraciones en Lean, un modelo necesita entender qué hay en la biblioteca, cómo se nombran los teoremas, qué tipos de tactic se usan con frecuencia y cómo interpretar los objetivos fallidos. Esto se parece mucho al autocompletado de código común, pero con restricciones más estrictas.

\begin{knowledgebox}{Semejanzas y diferencias entre Lean y el código}
La semejanza está en que ambos tienen sintaxis, tipos, bibliotecas, mensajes de error y retroalimentación ejecutable; la diferencia está en que el objetivo de Lean no es ejecutar un programa, sino demostrar proposiciones. Un programa puede apoyarse en pruebas que cubran los casos comunes; una demostración, en cambio, debe cumplirse en todos los casos.
\end{knowledgebox}

\subsection{El espacio de búsqueda de Proof Search}

La búsqueda de demostraciones no es una generación lineal. Un teorema puede requerir construir primero un lemma intermedio, elegir la variable de inducción, transformar la forma del objetivo, invocar teoremas existentes o demostrar primero que no existe un contraejemplo. El espacio de búsqueda es enorme y los caminos erróneos son muchos. La ventaja de la IA es que puede hacer intentos en paralelo y aprovechar conocimiento semántico previo; su desventaja es que tiende a generar pasos que parecen razonables pero que no están disponibles en la biblioteca.

\begin{warningbox}{Fallas comunes de la búsqueda de demostraciones}
El modelo puede conocer la idea matemática pero no saber cómo se llama el teorema correspondiente en la biblioteca de Lean; también puede conocer la tactic pero convertir el objetivo en una forma más difícil; o incluso puede generar una demostración que es válida en lenguaje natural pero a la que le faltan premisas en el sistema formal. La propia información de las fallas es, en sí misma, datos de entrenamiento.
\end{warningbox}

\subsection{Resumen de la sección}

Lo esencial de AI for Math no es «un modelo un poco más inteligente», sino descomponer la investigación matemática en un ciclo cerrado con retroalimentación: la intuición propone la dirección, la formalización restringe el proceso, el verificador proporciona retroalimentación estricta y el matemático juzga el significado.
